title: Conformal Mesh
alias: conformal, effective distance
todo:
------------------------------------

We approximate that the mesh sits on a flat two-dimensional *chart*: each node has a position $\vec{x}_i$, and the chart distance between two nodes is the ordinary Euclidean distance $d_{ij} = |\vec{x}_i - \vec{x}_j|$. The chart distance alone, however, treats every kilometre as equal, whereas a mosquito crossing a kilometre of continuous settlement (with houses, water containers and blood meals all the way) faces a very different journey to one crossing a kilometre of open fields between two villages.

Rather than add a new kind of object to the [[mesh]] to represent the land between villages, we absorb this into the distance itself, by giving the chart a *conformal* metric:
\begin{equation}
	{\rm d}s = e^{u(\vec{x})} |{\rm d}\vec{x}|
\end{equation}
where $u(\vec{x}) \geq 0$ is a dimensionless measure of how hard it is to travel through the space around $\vec{x}$. Where $u = 0$, distances are their chart values; where $u > 0$, they are stretched. A conformal metric changes lengths but not angles, and so it is the simplest modification which leaves the local geometry (``which neighbour is to the north'') untouched.

\section{The node approximation}

	We know $u$ only at the nodes, $u_i = u(\vec{x}_i)$, as a per-node [[mesh|NodeField]]. Assuming the path between two nodes is a straight line in the chart, along which $u$ varies slowly, the effective distance is the integral of the metric along it, which we approximate by the geometric mean of the conformal factors at the two ends:
	\begin{equation}
		d^{\rm eff}_{ij} = \int_{\vec{x}_i}^{\vec{x}_j} e^{u} \, |{\rm d}\vec{x}| \approx d_{ij} \exp\left(\frac{u_i + u_j}{2}\right)
	\end{equation}
	This has the properties we need:

	* It is symmetric, $d^{\rm eff}_{ij} = d^{\rm eff}_{ji}$.
	* It reduces to the chart distance between two urban nodes ($u_i = u_j = 0$).
	* It depends only on per-node quantities, so it needs no extra data beyond what the mesh already holds.

	As an example, with urban nodes assigned $u = 0$ and rural nodes $u = 1$, two neighbouring nodes a distance $d$ apart are effectively $d$ apart if both are urban, $e^{1/2}d \approx 1.65d$ if one is rural, and $e\,d \approx 2.72d$ if both are rural.

\section{How good is the approximation?}

	If $u$ varies linearly along the path, from $u_i$ to $u_j$, the integral can be done exactly:
	\begin{equation}
		\frac{d^{\rm eff, exact}_{ij}}{d^{\rm eff, approx}_{ij}} = \frac{\sinh(\Delta / 2)}{\Delta / 2}, \qquad \Delta = u_j - u_i
	\end{equation}
	which is always at least one, so the approximation slightly underestimates the distance. The error is small for moderate contrasts: about $4\%$ for $\Delta = 1$ (urban to rural, in the example above), and about $18\%$ for $\Delta = 2$.

	There is a second, opposing, error: the shortest path in a conformal metric is not the straight line, but bends around regions of high $u$ (a mosquito may prefer to follow a ribbon of settlement, rather than cross the fields directly). The straight line therefore overestimates the true effective distance. For a mesh with a fine resolution compared to the scale on which $u$ varies, both errors are small.

\section{Use in the code}

	The effective distance is computed in exactly one place, a single function of two node indices, which is called only when the [[environment]] is built: the dispersal probabilities between every connected pair of nodes are then calculated once from it, and cached. The rest of the code never sees a distance at all, only the cached probabilities, and so the cost of the metric (however elaborate) is paid only once.

	This also means the approximation above can be improved later, without any other change to the code. Possible refinements are:

	* A true line integral along the path, using a $u$ field sampled more finely than the nodes (for example, from a land-cover map).
	* Treating some regions as having very large $u$.

	The latter would be the natural way to represent barriers such as open water. However, whether water really is a barrier is an open question, since mosquitoes also travel by hitching a ride on human transport (boats crossing Lake Victoria, for example). Human-mediated movement is therefore planned as a separate set of edges (see [[mesh]]) which does not use this metric at all: water is never simply assumed to be impassable.

\section{Open questions}

	* How should $u$ be assigned in practice? From settlement density, from land cover, or simply from the urban/rural classification of the source data?
	* Should the dispersal kernel be applied to $d^{\rm eff}$ in the same functional form as it is to the chart distance, or does stretching the distance also change the shape of the kernel?
